%! TeX root = master.tex
% !TEX program = lualatex

\lecture{4}{2026-09-30}{The Upper Confidence Bound Algorithm}{}

\section{Recap: Bandits, Regret and Explore-then-Commit}

\subsection{Multi-armed bandits}

\begin{itemize}
    \item \textbf{Setup}: $k$ arms, each arm $i$ generates random rewards with distribution $P_i$ and mean $\mu_i$. The best arm has mean $\mu^*$. The reward distributions are \emph{unknown} to the learner at the beginning.
    \item \textbf{Interaction}: at each time $t = 1, 2, \dots, n$, the learner chooses an action $A_t \in \{1, \dots, k\}$ and observes a reward $X_t \sim P_{A_t}$. This is a sequential decision-making process.
    \item \textbf{Goal}: achieve \emph{sublinear regret}, i.e., eventually pull the best arm (almost) every time.
\end{itemize}

\subsection{Regret}

\begin{itemize}
    \item The \textbf{instantaneous regret} at round $t$ is $r_t = \mu^* - \mu_{A_t} = \Delta_{A_t}$. It measures the (expected) suboptimality of the arm played: a bad arm has a large regret, a good arm a small one, and the best arm has zero regret.
    \item The \textbf{cumulative regret} is $R_n = \sum_{t=1}^n \mathbb{E}[r_t]$.
    \item \textbf{Regret decomposition}:
    \[
    R_n = \sum_{t=1}^n \mathbb{E}[r_t] = \sum_{a \in \mathcal{A}} \Delta_a\, \mathbb{E}[T_a(n)].
    \]
\end{itemize}

\begin{remark}
Sublinear cumulative regret $\Leftrightarrow$ vanishing instantaneous regret. To minimize regret, bad arms must be eliminated quickly, and eventually the best arm(s) must be identified and pulled. But, because of randomness, good and bad arms can \emph{never} be identified with certainty!
\end{remark}

\subsection{Explore-then-Commit (ETC)}

\begin{itemize}
    \item Explore until some fixed time $mk$ ($m$ pulls per arm). The exploration phase gives an estimate $\hat{\mu}_i$ of $\mu_i$ for every $i$.
    \item If $\hat{\mu}_i \approx \mu_i$, then $i^* = \arg\max_i \hat{\mu}_i$ identifies the best arm. Commit to (keep pulling) $i^*$ for rounds $mk+1$ to $n$.
\end{itemize}

\textbf{Intuition.} If the best arm is correctly identified, there is no regret in the commit phase, hence $R_n \propto m$ (great!). If a suboptimal arm is selected, a constant positive regret is paid at \emph{every} round of the commit phase, hence $R_n \propto n$ (bad!).

\textbf{Tradeoff.} A very small $m$ means $|\hat{\mu}_i - \mu_i|$ is large, so we are likely to choose a suboptimal arm (large regret). A very large $m$ means excessive exploration (also large regret). Tuning $m$ optimally requires knowing the gap $\Delta$ in advance.

\section{Adaptive Explore-then-Commit}

\subsection{Key idea: $m$ as a stopping time}

\begin{itemize}
    \item \textbf{Key idea}: think of $m$ as a \emph{stopping time}. It is random and can depend on past data; it signals when to stop exploring and start exploiting. Ideally, we stop when we are sufficiently confident about which arm(s) are optimal.
    \item With a \emph{fixed} $m$, two things can go wrong: we might not be sure about the optimal arm at time $mk$ (\emph{under-exploration}), or we might have been very sure about the optimal arm much earlier (\emph{over-exploration}).
    \item \textbf{Idea}: decide when to stop by looking at the data. Keep pulling all arms equally and computing the $\hat{\mu}_i$, until we are (reasonably) certain that one arm is better than all the others.
\end{itemize}

How do we quantify certainty? With \textbf{confidence intervals}.

\subsection{Confidence intervals: recap}

Let $X_1, \dots, X_n$ be i.i.d.\ samples of a random variable with unknown mean $\mu$, and let $\hat{\mu}_n = \frac{1}{n}\sum_i X_i$ be the estimate of $\mu$. Typically, $\hat{\mu}_n$ is close to $\mu$, but not exactly equal to it. We can say:
\[
|\hat{\mu}_n - \mu| \leq \epsilon \quad \text{with probability at least } 1 - \delta,
\]
or equivalently, $\mu \in [\hat{\mu}_n - \epsilon, \hat{\mu}_n + \epsilon]$ with high probability.

\begin{center}
\fbox{Inserer ici slide 6}
\end{center}

\textbf{Confidence interval $\epsilon$ vs.\ confidence level $\delta$.} Suppose the $X_i$ are $\sigma$-subgaussian. The subgaussian concentration inequality (lecture 3) states
\[
\mathbb{P}\big(|\hat{\mu}_n - \mu| \geq \epsilon\big) \leq 2\exp\left(-\frac{n\epsilon^2}{2\sigma^2}\right) =: \delta.
\]
Given $\epsilon$ we can compute $\delta$; alternatively, we fix $\delta$ and compute $\epsilon$:
\[
\epsilon = \sigma\sqrt{\frac{2\ln(2/\delta)}{n}}.
\]
\begin{itemize}
    \item For fixed $\delta$, $\epsilon$ decreases as $n$ increases (more data, narrower interval).
    \item For fixed $n$, $\epsilon$ decreases as $\delta$ increases (accepting a larger error probability gives a narrower interval).
\end{itemize}

\textbf{Upper and lower confidence bounds.} With $n$ the number of samples of the arm considered,
\[
\mathrm{UCB}(n,\delta) = \hat{\mu}_n + \sigma\sqrt{\frac{2\log(2/\delta)}{n}}, \qquad
\mathrm{LCB}(n,\delta) = \hat{\mu}_n - \sigma\sqrt{\frac{2\log(2/\delta)}{n}}.
\]

\subsection{Using confidence intervals to decide when to commit}

How can we confidently decide that one arm is better than another? Pretend that each arm's mean lies in its confidence interval with certainty.
\begin{itemize}
    \item If the confidence intervals \emph{overlap}, the best arm is undecided: keep exploring.
    \item If the confidence intervals are \emph{disjoint}, the best arm is clear: commit.
\end{itemize}
\textbf{Strategy}: keep exploring until the confidence intervals separate.

\begin{center}
\fbox{Inserer ici slide 7}
\end{center}

\textbf{Data-dependent exploration time.} With sufficient exploration, the confidence intervals eventually separate.
\begin{itemize}
    \item When $\mu_1 - \mu_2$ is \emph{large}, the intervals separate early, so $m$ tends to be small ($\mathbb{E}[m]$ small).
    \item When $\mu_1 - \mu_2$ is \emph{small}, the intervals separate only after a long time, so $m$ is large ($\mathbb{E}[m]$ large).
\end{itemize}
The exploration time is thus data-dependent, not fixed a priori. (Slide example: $m = 20$ is enough to separate the intervals for a large $\Delta$; for the same $m=20$ the intervals still overlap for a small $\Delta$, and one needs e.g.\ $m = 100$ to separate them.)

\begin{center}
\fbox{Inserer ici slide 8}
\end{center}

\subsection{Pseudo-code}

\begin{definition}[Adaptive ETC (ETC with confidence intervals)]
\textbf{Initialize}: $T_a(0) = 0$, $\hat{\mu}_a(0) = 0$ for all arms $a \in \mathcal{A}$; $\texttt{committed\_arm} = \texttt{None}$.\\
\textbf{Hyperparameters}: $\delta \in (0,1)$ and a minimum number of pulls $m_0$.

\medskip
\noindent For $t = 1, 2, \dots, n$:
\begin{itemize}
    \item If $\texttt{committed\_arm} \neq \texttt{None}$: $A_t = \texttt{committed\_arm}$. Else: $A_t = (t \bmod k) + 1$ (round robin).
    \item \textbf{Check for commitment} if $t \geq m_0 k$:
    \begin{itemize}
        \item compute $\hat{\mu}_a(t)$ for all arms and $a^*(t) = \arg\max_a \hat{\mu}_a(t)$;
        \item compute $\mathrm{UCB}_a(t,\delta)$ and $\mathrm{LCB}_a(t,\delta)$ for all arms;
        \item if $\mathrm{LCB}_{a^*}(t,\delta) > \mathrm{UCB}_a(t,\delta)$ for all $a \neq a^*$, then $\texttt{committed\_arm} = a^*$.
    \end{itemize}
\end{itemize}
\end{definition}

\begin{remark}
\begin{itemize}
    \item The empirically best arm is committed to only when its \emph{lower} confidence bound exceeds the \emph{upper} confidence bound of \emph{every} other arm, i.e., when its confidence interval dominates all the others.
    \item The minimum number of pulls $m_0$ prevents committing on the basis of (nearly) no data.
    \item Heuristically, if all the confidence intervals contain the true means (which happens with high probability when $\delta$ is small) then $\mu_{a^*} \ge \mathrm{LCB}_{a^*} > \mathrm{UCB}_a \ge \mu_a$ for all $a \neq a^*$, so the commitment is correct.
\end{itemize}
\end{remark}

\subsection{Regret curves of adaptive ETC}

\begin{center}
\fbox{Inserer ici slide 11}
\end{center}

Salient features of adaptive ETC (simulations on a two-armed bandit):
\begin{itemize}
    \item the exploration period is \emph{random};
    \item the expected cumulative regret is \emph{sublinear}.
\end{itemize}

\begin{center}
\fbox{Inserer ici slide 12}
\end{center}

\begin{itemize}
    \item A larger $\Delta$ means the algorithm commits earlier.
    \item As for fixed ETC, the regret is highest for a \emph{moderate} $\Delta$ (e.g.\ $\Delta = 0.1$ in the simulations).
\end{itemize}

\begin{center}
\fbox{Inserer ici slide 13}
\end{center}

\textbf{Importance of the hyperparameter $\delta$:}
\begin{itemize}
    \item Very small $\delta$ $\Rightarrow$ too conservative (over-exploration).
    \item Very large $\delta$ $\Rightarrow$ too optimistic: there is a chance of a mistaken commitment (some runs commit to the wrong arm and have linear regret).
    \item \textbf{Rule of thumb}: the larger the horizon $n$, the smaller $\delta$ should be (mistakes are more costly).
\end{itemize}

\section{Beyond Two Arms: Sequential Elimination}

\subsection{ETC for $k > 2$ arms}

How can the ETC principle be extended to $k > 2$ arms? We need to pull each arm a few times to estimate its reward, so exploration of all arms is necessary. As in the two-armed case, with more exploration the confidence interval of each arm keeps shrinking, eventually the intervals become disjoint, and the best arm becomes clear. There are, however, two different interpretations:

\begin{itemize}
    \item \textbf{View 1} (identify the best arm): keep exploring until the best arm is identified; commit to arm $i$ once its confidence interval dominates the confidence intervals of all other arms.
    \item \textbf{View 2} (eliminate bad arms): keep eliminating arms that are clearly not the best; if the confidence interval of arm $i$ dominates the one of arm $j$, stop pulling arm $j$.
\end{itemize}
These two views are \emph{equivalent with two arms} but \emph{different} in the multi-armed case.

\subsection{From ETC to sequential elimination}

\begin{itemize}
    \item \textbf{Algorithm from View 1} (adaptive ETC): all arms are pulled equally often until the commit time $mk$. Likely behavior: the worse arms clearly separate out well before the competitive ones, but they keep being pulled $\Rightarrow$ wasteful exploration on bad arms, hence large regret.
    \item \textbf{Algorithm from View 2} (\textbf{Sequential Elimination}): initially all arms are pulled in round robin. The algorithm proceeds in \emph{phases}, eliminating one arm per phase; the length of each phase is random. The instantaneous regret decreases from phase to phase, and eventually only the best arm remains.
\end{itemize}

Sequential elimination is the ``right'' extension of ETC to multi-armed bandits (to be explored further in the homework).

\textbf{Comparison on a $4$-armed Bernoulli bandit with means $[0.2, 0.4, 0.6, 0.8]$.}
\begin{itemize}
    \item \emph{Regret curves}: adaptive ETC shows no commitment in some runs; for sequential elimination the regret slope decreases in phases.
    \item \emph{Confidence intervals}: in adaptive ETC all the intervals have the same length (all arms are pulled equally), and some overlap may remain; in sequential elimination, the better arms have smaller confidence intervals: they stay in the race longer, are pulled more often, and hence have \emph{narrower} intervals, while eliminated arms keep wide ones.
\end{itemize}

\begin{center}
\fbox{Inserer ici slide 16}
\end{center}
\begin{center}
\fbox{Inserer ici slide 17}
\end{center}

\section{The Upper Confidence Bound (UCB) Algorithm}

UCB is a very popular algorithm for multi-armed bandits. Like flexible ETC and sequential elimination it uses confidence bounds, but it leads to better regret bounds in practice.

\subsection{Definition}

\textbf{Notation.} $\hat{\mu}_a(t)$ is the estimate of the mean reward of arm $a$ after $t$ overall pulls and $T_a(t)$ is the number of pulls of arm $a$ at time $t$ (``at time $t$'' refers to after $t$ overall pulls). We assume \textbf{$1$-subgaussian} rewards.

\begin{definition}[Upper confidence bound]
The upper confidence bound of arm $a$ at time $t$ is
\[
\mathrm{UCB}_a(t,\delta) = \hat{\mu}_a(t) + \sqrt{\frac{2\log(1/\delta)}{T_a(t)}},
\]
where $\mathrm{UCB}_a(t,\delta) = +\infty$ initially (when $T_a(t) = 0$).
\end{definition}

\begin{definition}[UCB algorithm]
For $t = 1, 2, \dots, n$:
\begin{itemize}
    \item choose $A_t = \arg\max_a \mathrm{UCB}_a(t-1, \delta)$;
    \item observe $X_t$ and update the $\mathrm{UCB}_a(t,\delta)$.
\end{itemize}
\end{definition}

\begin{remark}
\begin{itemize}
    \item Because $\mathrm{UCB}_a = +\infty$ for an unplayed arm, the algorithm starts by pulling each arm once.
    \item Only the upper bound is needed (a one-sided bound), hence $\ln(1/\delta)$ rather than $\ln(2/\delta)$ as in the two-sided confidence intervals.
    \item The only hyperparameter is $\delta$. The theorem below uses $\delta = 1/n^2$.
\end{itemize}
\end{remark}

\subsection{Exploration-exploitation tradeoff in UCB}

The index of an arm decomposes as
\[
\mathrm{UCB}_a(t,\delta) = \underbrace{\hat{\mu}_a(t)}_{\text{empirical mean}} + \underbrace{\sqrt{\frac{2\log(1/\delta)}{T_a(t)}}}_{\text{exploration bonus}}.
\]
The algorithm pulls an arm more often if it is
\begin{itemize}
    \item either \textbf{promising}: the empirical mean $\hat{\mu}_a(t)$ is large (exploitation), or
    \item \textbf{not well explored}: $T_a(t)$ is small, so the exploration bonus is large (exploration).
\end{itemize}

\textbf{Why is this a reasonable strategy?}
\begin{itemize}
    \item With high probability, $\mathrm{UCB}_{a^*}(t,\delta) > \mu_{a^*}$ at all times.
    \item If a suboptimal arm $a \neq a^*$ is pulled, then $\mathrm{UCB}_a(t,\delta) > \mathrm{UCB}_{a^*}(t,\delta)$.
    \item With enough pulls, $\mathrm{UCB}_a(t,\delta) \approx \mu_a$ (the bonus vanishes), and $\mu_a < \mu_{a^*} < \mathrm{UCB}_{a^*}(t,\delta)$, so $a$ stops being chosen.
    \item Hence suboptimal arms are not pulled too many times!
\end{itemize}

\begin{center}
\fbox{Inserer ici slide 19}
\end{center}

\subsection{UCB vs.\ ETC: optimism in the face of uncertainty}

\begin{itemize}
    \item UCB follows the \textbf{Optimism in the Face of Uncertainty (OFU)} principle: it plays assuming that arm rewards are the highest they can \emph{reasonably} be, i.e., it pretends that $\mu_a = \mathrm{UCB}_a(t,\delta)$.
    \item ETC (and sequential elimination) is more \emph{conservative}: it takes a worst-case approach, where arm rewards could be anything within the respective confidence intervals. It only commits to the best arm if it is ``sure'' that it is the best (its interval must dominate all the others), and only eliminates a bad arm if it is ``sure'' it is not the best.
    \item UCB has a \textbf{soft} transition from exploration to exploitation.
    \item ETC has a \textbf{hard} transition (explore phase vs.\ commit phase), which makes it difficult to correct a mistake.
\end{itemize}

\begin{center}
\fbox{Inserer ici slide 20}
\end{center}

\subsection{Regret performance of UCB (simulations)}

\begin{center}
\fbox{Inserer ici slide 21}
\end{center}
\begin{center}
\fbox{Inserer ici slide 22}
\end{center}

Points to note:
\begin{itemize}
    \item For each sample path, the regret growth decays gently with time (contrast with ETC, where the slope drops abruptly at commitment, or never drops if the wrong arm is chosen).
    \item Exploration never stops completely.
    \item The mean regret is sublinear. Is it $O(\sqrt{n})$? $O(\log n)$? (The theorem below answers: $O(\log n)$ in the instance-dependent sense.)
\end{itemize}

\section{Regret Bound for UCB}

\begin{theorem}[Regret of UCB]
The UCB algorithm applied to a stochastic $1$-subgaussian bandit with $k$ arms and an error tolerance $\delta = 1/n^2$ satisfies
\[
R_n \leq 3\sum_{i=1}^k \Delta_i + \sum_{i:\,\Delta_i > 0} \frac{16\ln(n)}{\Delta_i}.
\]
\end{theorem}

\begin{remark}
\begin{itemize}
    \item The first term, $3\sum_i \Delta_i$, is a constant independent of $n$; the second term is \emph{logarithmic} in $n$. Hence UCB achieves \textbf{logarithmic} (and in particular sublinear) regret.
    \item The bound is \emph{instance-dependent}: it depends on the gaps $\Delta_i$. A smaller gap makes an arm harder to distinguish from the best arm and increases the regret bound ($\propto 1/\Delta_i$). Note that this bound alone is not informative when some $\Delta_i \to 0$.
    \item \emph{(Side remark, not on the slides.)} Since the regret of an arm with gap $\Delta_i < \Delta_0$ is at most $n \Delta_0$ in total, splitting the arms according to their gap and optimizing $\Delta_0$ gives a worst-case (problem-independent) bound of order $O(\sqrt{kn\ln n})$.
    \item Compared with ETC, UCB needs to know neither $\Delta$ nor a tuned $m$.
\end{itemize}
\end{remark}

\subsection{Notation: $\hat{\mu}_i(t)$ vs.\ $\hat{\mu}_{is}$}

\begin{itemize}
    \item $\hat{\mu}_i(t)$ is the estimator of arm $i$ \emph{at time $t$}.
    \item $\hat{\mu}_{is}$ is the estimator \emph{after $s$ pulls of arm $i$}, i.e., the average of $s$ i.i.d.\ samples (the first $s$ rewards generated by arm $i$).
    \item The two are related by $\hat{\mu}_i(t) = \hat{\mu}_{i T_i(t)}$. Remember that $T_i(t) \leq t$ for all $t$.
\end{itemize}
The key point is that $\hat{\mu}_{is}$ for a \emph{fixed, deterministic} $s$ is an average of $s$ i.i.d.\ samples, so the concentration results of lecture 3 apply to it directly, whereas $\hat{\mu}_i(t)$ involves the random number of pulls $T_i(t)$.

\begin{center}
\fbox{Inserer ici slide 25}
\end{center}

\subsection{Proof outline}

\begin{center}
\fbox{Inserer ici slide 24}
\end{center}

Without loss of generality, assume that arm $1$ is optimal ($\mu_1 = \mu^*$) and fix a suboptimal arm $i$ ($\Delta_i > 0$). By the regret decomposition, it suffices to bound $\mathbb{E}[T_i(n)]$. The proof has four steps:
\begin{enumerate}
    \item Define a ``good event'' $G_i$ such that: (1) $\mu_1$ is \emph{always} (for all $t \in [n]$) below its own UCB; and (2) $\mu_1$ is above the UCB of arm $i$ after $u_i$ pulls of arm $i$.
    \item Show that if $G_i$ holds, then $T_i(n) \leq u_i$.
    \item Show that if $u_i$ is large enough, then $G_i$ is very likely.
    \item Make a specific choice for $u_i$ that allows us to bound $\mathbb{E}[T_i(n)]$ (without conditioning on $G_i$).
\end{enumerate}

\subsection{Step 1: the good event $G_i$}

Let $u_i \in \mathbb{N}$ be a (deterministic) number to be chosen later, and define
\[
G_i = \left\{\mu_1 < \min_{t \in [n]} \mathrm{UCB}_1(t,\delta)\right\} \cap \left\{\hat{\mu}_{i u_i} + \sqrt{\frac{2}{u_i}\ln\frac{1}{\delta}} < \mu_1\right\}.
\]
\begin{itemize}
    \item First event: $\mu_1$ is always inside the confidence interval of arm $1$ (more precisely, below its UCB).
    \item Second event: at a time $t$ where $T_i(t) = u_i$, the UCB of arm $i$ does not exceed $\mu_1$.
\end{itemize}

Decomposing the expectation according to $G_i$ and using $T_i(n) \leq n$ on $G_i^c$,
\[
\mathbb{E}[T_i(n)] = \mathbb{E}\big[\mathbb{I}\{G_i\}\, T_i(n)\big] + \mathbb{E}\big[\mathbb{I}\{G_i^c\}\, T_i(n)\big] \leq u_i + \mathbb{P}(G_i^c)\, n.
\]
There is a tradeoff in the choice of $u_i$:
\begin{itemize}
    \item $u_i$ too large $\Rightarrow$ the regret bound is too large;
    \item $u_i$ too small $\Rightarrow$ $\mathbb{P}(G_i^c)$ is too large $\Rightarrow$ the regret bound is too large.
\end{itemize}

\begin{center}
\fbox{Inserer ici slide 26}
\end{center}

\subsection{Step 2: if $G_i$ holds, then $T_i(n) \leq u_i$}

\begin{proof}
By contradiction: assume that $G_i$ holds and $T_i(n) > u_i$. Then arm $i$ was pulled more than $u_i$ times before $n$, so there is a round $t$ at which arm $i$ was pulled for the $(u_i+1)$-th time, i.e., $T_i(t-1) = u_i$ and $A_t = i$. At this round,
\begin{align*}
\mathrm{UCB}_i(t-1,\delta) &= \hat{\mu}_i(t-1) + \sqrt{\frac{2\ln(1/\delta)}{T_i(t-1)}} \\
&= \hat{\mu}_{i u_i} + \sqrt{\frac{2}{u_i}\ln\frac{1}{\delta}} \qquad (T_i(t-1) = u_i) \\
&< \mu_1 \qquad \text{(second event of } G_i) \\
&< \mathrm{UCB}_1(t-1,\delta) \qquad \text{(first event of } G_i).
\end{align*}
Hence arm $1$ has a strictly larger UCB than arm $i$ at time $t-1$, so the algorithm could not have chosen arm $i$ at time $t$: contradiction!
\end{proof}

\begin{center}
\fbox{Inserer ici slide 27}
\end{center}

\subsection{Step 3: $G_i^c$ is unlikely}

By De Morgan,
\[
G_i^c = \left\{\mu_1 \geq \min_{t \in [n]} \mathrm{UCB}_1(t,\delta)\right\} \cup \left\{\hat{\mu}_{i u_i} + \sqrt{\frac{2}{u_i}\ln\frac{1}{\delta}} \geq \mu_1\right\}.
\]
We bound the probability of each of the two events.

\textbf{First event.} Since $\mathrm{UCB}_1(t,\delta) = \hat{\mu}_{1 T_1(t)} + \sqrt{2\ln(1/\delta)/T_1(t)}$ and $T_1(t) \in [n]$, the minimum over the $n$ rounds $t \in [n]$ is at least the minimum over all possible numbers of pulls $s \in [n]$ (the latter ranges over at least as many values, which can only lower the minimum):
\[
\left\{\mu_1 \geq \min_{t \in [n]} \mathrm{UCB}_1(t,\delta)\right\} \subset \left\{\mu_1 \geq \min_{s \in [n]} \left(\hat{\mu}_{1s} + \sqrt{\frac{2}{s}\ln\frac{1}{\delta}}\right)\right\} = \bigcup_{s \in [n]} \left\{\mu_1 \geq \hat{\mu}_{1s} + \sqrt{\frac{2}{s}\ln\frac{1}{\delta}}\right\}.
\]
Each $\hat{\mu}_{1s}$ is an average of $s$ i.i.d.\ samples, so by the union bound and the subgaussian concentration inequality (with $\epsilon = \sqrt{2\ln(1/\delta)/s}$, so that $\exp(-s\epsilon^2/2) = \delta$),
\[
\mathbb{P}\left(\mu_1 \geq \min_{t \in [n]} \mathrm{UCB}_1(t,\delta)\right) \leq \sum_{s=1}^n \mathbb{P}\left(\mu_1 \geq \hat{\mu}_{1s} + \sqrt{\frac{2}{s}\ln\frac{1}{\delta}}\right) \leq n\delta.
\]

\textbf{Second event.} Here we need an assumption on $u_i$: assume that $u_i$ is large enough so that
\[
\Delta_i - \sqrt{\frac{2\ln(1/\delta)}{u_i}} \geq c\Delta_i \qquad \text{for some constant } 0 < c < 1.
\]
Using $\mu_1 = \mu_i + \Delta_i$,
\begin{align*}
\mathbb{P}\left(\hat{\mu}_{i u_i} + \sqrt{\frac{2}{u_i}\ln\frac{1}{\delta}} \geq \mu_1\right) &= \mathbb{P}\left(\hat{\mu}_{i u_i} - \mu_i \geq \Delta_i - \sqrt{\frac{2}{u_i}\ln\frac{1}{\delta}}\right) \\
&\leq \mathbb{P}\big(\hat{\mu}_{i u_i} - \mu_i \geq c\Delta_i\big) \\
&\leq \exp\left(-\frac{u_i c^2\Delta_i^2}{2}\right),
\end{align*}
where the last step is the Chernoff bound for subgaussian averages ($u_i$ is the number of terms in $\hat{\mu}_{i u_i}$).

\textbf{Combining} the two bounds (union bound):
\[
\mathbb{P}(G_i^c) \leq n\delta + \exp\left(-\frac{u_i c^2\Delta_i^2}{2}\right).
\]

\begin{center}
\fbox{Inserer ici slide 28}
\end{center}
\begin{center}
\fbox{Inserer ici slide 29}
\end{center}

\subsection{Step 4: choice of $u_i$ and conclusion}

Back to the bound on the expected number of pulls:
\[
\mathbb{E}[T_i(n)] \leq u_i + \mathbb{P}(G_i^c)\, n \leq u_i + n\left(n\delta + \exp\left(-\frac{u_i c^2\Delta_i^2}{2}\right)\right).
\]
We set $u_i$ as small as possible given the constraint $\Delta_i - \sqrt{2\ln(1/\delta)/u_i} \geq c\Delta_i$, i.e.,
\[
u_i = \left\lceil \frac{2\ln(1/\delta)}{(1-c)^2\Delta_i^2} \right\rceil.
\]
(There is no guarantee that this is $\leq n$, but if it is not, the bound on $\mathbb{E}[T_i(n)]$ holds trivially since $T_i(n) \leq n$.)

Recall $\delta = 1/n^2$, so $\ln(1/\delta) = 2\ln n$ and $n^2\delta = 1$. The exponential term becomes
\[
\exp\left(-\frac{u_i c^2\Delta_i^2}{2}\right) \leq \exp\left(-\frac{c^2\ln(1/\delta)}{(1-c)^2}\right) = \exp\left(-\frac{2c^2\ln n}{(1-c)^2}\right) = n^{-\frac{2c^2}{(1-c)^2}}.
\]
Therefore
\[
\mathbb{E}[T_i(n)] \leq \frac{2\ln(1/\delta)}{(1-c)^2\Delta_i^2} + 1 + 1 + n^{1 - \frac{2c^2}{(1-c)^2}},
\]
where one ``$+1$'' comes from the ceiling in $u_i$ and the other from $n^2\delta = 1$. We choose $c$ such that the last term does not blow up, e.g., $c = 0.5$: then $\frac{2c^2}{(1-c)^2} = 2$, so $n^{1-2} = 1/n \leq 1$ and $\frac{2\ln(1/\delta)}{(1-c)^2} = \frac{4\ln n}{1/4} = 16\ln n$. Hence
\[
\mathbb{E}[T_i(n)] \leq 3 + \frac{16\ln n}{\Delta_i^2}.
\]
Finally, by the regret decomposition,
\[
R_n = \sum_{i:\,\Delta_i > 0} \Delta_i\,\mathbb{E}[T_i(n)] \leq 3\sum_{i=1}^k \Delta_i + \sum_{i:\,\Delta_i > 0}\frac{16\ln n}{\Delta_i}. \qquad \blacksquare
\]

\begin{center}
\fbox{Inserer ici slide 30}
\end{center}

\subsection{UCB regret plots}

Simulations with Bernoulli rewards of means $[0.5 - \Delta/2,\; 0.5 + \Delta/2]$:
\begin{itemize}
    \item The regret is sublinear for all values of $\Delta$, given a fixed hyperparameter $\delta$.
    \item The cumulative regret increases inversely with $\Delta$ (a smaller gap is harder).
    \item The slope gradually reduces with time, signifying the shift from exploration to exploitation.
    \item The initial regret slope is higher for a large $\Delta$, due to the more expensive exploration (each pull of the bad arm costs $\Delta$).
\end{itemize}

\begin{center}
\fbox{Inserer ici slide 31}
\end{center}

\textbf{UCB vs.\ ETC} (figure from Chapter 7 of Lattimore \& Szepesv\'ari): UCB performs similarly to ETC, \emph{without knowing the optimal $m$} (i.e., without knowing the suboptimality gap).

\begin{center}
\fbox{Inserer ici slide 32}
\end{center}

\section{Summary}

\begin{itemize}
    \item \textbf{Explore-then-Commit with a flexible stopping time}: the exploration phase is chosen adaptively, based on data (commit when the confidence interval of the empirically best arm dominates all the others); it gives sublinear regret, and the hyperparameter $\delta$ trades off over-exploration (too small) against mistaken commitment (too large). It extends to \textbf{sequential elimination} for multiple arms.
    \item \textbf{Main takeaway: the UCB algorithm}:
    \begin{itemize}
        \item quantifies uncertainty with confidence intervals: $\mathrm{UCB}_a(t,\delta) = \hat{\mu}_a(t) + \sqrt{2\log(1/\delta)/T_a(t)}$;
        \item acts \emph{optimistically in the face of uncertainty} (always pulls the arm with the largest UCB);
        \item smooth transition from exploration to exploitation (no hard commitment, no knowledge of $\Delta$ needed);
        \item regret bound (with $\delta = 1/n^2$, $1$-subgaussian rewards): $R_n \leq 3\sum_i \Delta_i + \sum_{i:\Delta_i>0} 16\ln(n)/\Delta_i$.
    \end{itemize}
    \item Proof technique: good event $G_i$ (the optimal arm's mean always below its UCB, and arm $i$'s UCB below $\mu_1$ after $u_i$ pulls) $\Rightarrow$ $T_i(n) \leq u_i$; bound $\mathbb{P}(G_i^c)$ by a union bound and Chernoff; choose $u_i \propto \ln(n)/\Delta_i^2$.
    \item Reading: Lattimore \& Szepesv\'ari, Chapter 7 (especially Theorem 7.1).
\end{itemize}
